\summarychapter{42}{Graffiti on a Wall of Shame}{Twentieth-Century Ideologies}

\begin{inbrief}
In the twentieth century the Christian message had to compete with political creeds that behaved
like religions: nationalism, communism and individualism. Out of the First World War came two
totalitarian systems, Nazism on the right and Soviet communism on the left, and both attacked the
churches. The chapter recounts how Christians fared under Hitler and under Stalin, what the Second
World War did to Christianity, and how the Cold War divided the world.
\question{How has Christianity, with its message of peace and freedom, endured through these years of
international hostilities?}
\end{inbrief}

In the summer of 1961 the East German police, urged on by the Soviet premier Nikita Khrushchev, tried
to stop the flight of East Germans to West Berlin. When tighter travel rules failed, they built
twenty-six miles of barbed wire and concrete across the city like an ugly scar. Berliners called it
the \emph{Schandmauer}, the wall of shame. East of it stood the somber buildings of the Communist
world with its gospel of a future classless utopia; west of it the shops and cinemas of the ``free
West,'' seeking wealth and happiness now; and alongside lay unmarked bunkers, silent memorials of
Hitler's Third Reich. Across such bloodstained walls the Christian message has been scribbled like
graffiti among other messages just as arresting and perhaps more demanding: Nazism, Marxism,
capitalism.

\ssection{Post-Christian Gods for the Masses}
The historian Arnold Toynbee suggested that the twentieth century saw the great world religions
displaced by three post-Christian ideologies: nationalism, communism and individualism. Each takes on
the character of a religion. Each makes an ultimate demand (patriotism, class struggle, secular
humanism) and has its sacred symbols and ceremonies, inspired writings, dogmas, saints and
charismatic leaders. Shelley calls the theory only partly right. Westerners have reason to read the
century as a process of secularization, in which the ideologies displaced or tamed the religions.
Many predicted that it would continue: reason and science would explain more and more, educated
people would need religion less, fewer would believe in God, and the God they believed in would be
less of a God. For the world as a whole, however, the forecast is ``ridiculously inaccurate.'' The
world is more religious than ever, and Western believers who think Christianity in retreat may be
unaware of its unparalleled expansion. Before the century ended theorists had reversed much of
Toynbee's view, concluding that the future is better understood through religious and cultural
issues than through political and economic ones.

Totalitarian governments did arise, usually led by a dictator or a small elite in control of armed
force. With sophisticated psychological methods the rulers turn the minds and feelings of the people
against the regime's enemies. Propaganda, control of the media and regulation of the economy are
meant to produce a new kind of person with no hunger for personal freedom.

The roots lie in the First World War. The new religion of early twentieth-century Europe was
nationalism. Pan-Germanism and Pan-Slavism set the powers against each other in the Balkans, and
militarism, economic imperialism and power politics made a tinderbox. The spark came on June 28,
1914, when a student inspired by Serbian nationalism assassinated the heir to the throne of
Austria-Hungary. By August Germany and Austria, the Central Powers, faced France, Russia and Britain,
the Allies, and in the end twenty-seven nations were drawn in. For the first time the world knew
``total war,'' in which civilians gave up liberties and possessions, and each side believed it was
on a righteous crusade. When the United States entered on April 6, 1917, President Woodrow Wilson,
a Presbyterian minister's son, called it a war ``to make the world safe for democracy.'' While
America mobilized, Russia collapsed. In November 1917 the Bolsheviks, the militant, self-appointed
vanguard of a new socialist society, seized power, and early in 1918 they made peace with Germany.
A million American soldiers enabled the Allies to bring Germany to an armistice and to the Paris
Peace Conference. The terms were a strange blend of naive idealism and revenge. Wilson obtained a
League of Nations, which his own people refused to support, and the other Allies extracted a
confession of guilt and crippled Germany for years. The Germans never forgot, and Nazism grew out of
that vindictiveness.

\ssection{The Rise of Nazism}
The Nazis taught the world what totalitarianism means. Theirs was dictatorship of the right, or
fascism. Such governments answer personal frustration and social and economic tension by stressing
class unity and reaffirming traditional values. They glorify the nation, defined by its mission, its
race or the state itself, and they permit private property and capitalist enterprise under tight
control. Right-wing regimes sprang up all over Europe after the war, but the greatest was German
National Socialism. The Protestant churches of Luther's land lost millions to this political
religion; many were compromised or coopted, and others abandoned traditional Christianity. Most
Protestant clergy were monarchists with no sympathy for socialism or democracy, so industrial
workers came to hate a reactionary church, while the cultured classes had turned to a romantic,
heroic, aristocratic and often pantheistic view of the German past.

The leader was the Austrian-born Adolf Hitler, named chancellor on January 30, 1933; within two years
he had complete control of the government. The Nazis believed in the absolute unity of the German
people under the leader (\emph{der F\"uhrer}) and in applying this leadership principle to every
structure of the nation, so as to create an ideal super-community. They dropped Enlightenment goals
such as individual freedom and tolerance but kept the Enlightenment method of criticizing the Bible's
traditions by reason. The result was an invented Christianity whose ``Nazi Jesus'' despised the Jews
and their testament; Christians, they said, had mistakenly set him in a Jewish framework. As is
typical of Enlightenment thinkers, Shelley remarks, they only pretended to criticize tradition. In
fact they exchanged the tradition of Jesus and the Jewish people for another, the idealized primitive
past of Wagner's operas and the Germanic sagas. Race was central. ``Soil'' and ``blood'' set Germans
apart from all other peoples, and foreign ideas and persons, above all Jewish ones, were corrupting.

From this came a barbaric anti-Semitism. To recover its lost innocence Germany must purge itself, and
the Jews were the scapegoat: the ``culture-destroying race'' that had given the world both capitalism
and Marxism. Hitler declared that Christianity itself was a Jewish plot, ``the heaviest blow that
ever struck humanity,'' and Bolshevism its illegitimate child. The Nazis first stripped German Jews
of citizenship and pressed them to emigrate. When their armies overran eastern Europe, with its much
larger Jewish population, death squads killed thousands on the spot. The concentration camps had
been created to break the regime's opponents; from 1940 a new kind of camp was designed for the
``final solution,'' the extermination of all the Jews of Europe. The largest was Auschwitz in Poland.
Men, women and children were transported to these ``death factories'' and starved, shot, gassed or
used for medical experiments. Reasonable estimates put the Jewish dead of the Holocaust at six
million.

\ssection{Christians under Hitler}
The dedicated Christian was no more secure than the Jew. Hitler, born and raised a Catholic, had
abandoned whatever Christian principles he once held for his faith in Germany's regeneration through
National Socialism. On his way to power he courted Christian support by appealing to national pride
and pretending to favor the churches' place in the state. Churchmen, like millions of Germans, had
felt the shock of defeat and dreamed of a national rebirth. Catholics, strong in the republic of the
1920s, endorsed the new government and supported the concordat that Hitler signed with the pope in
1933, which guaranteed freedom to practice the Catholic religion. It raised his prestige greatly and
removed Catholics from German politics, and he had no intention of keeping his side of it.

Among Protestants the German Christians sought closer ties with the Nazis. They wanted to unite the
twenty-eight regional churches under one bishop and elected to the post Ludwig M\"uller, a fervent
Nazi. They brought the leadership principle into church government and adopted the ``Aryan
paragraph,'' which dismissed everyone of Jewish origin from church positions. In 1933 they claimed
three thousand of the seventeen thousand Protestant pastors. Against them a group led by Martin
Niem\"oller formed the Pastors' Emergency League and set up an alternative church government, the
Confessing Church, at its height probably about as large as its rival. Between the two stood the
timid majority, who in the end obeyed Hitler without open protest.

In May 1934 the Confessing Church stated its convictions in the Barmen Declaration, written largely
by Karl Barth. It called the German churches back to the central truths of Christianity, rejected the
totalitarian claims of the state, and bound the church to revelation, above all in Jesus. It ruled
out natural theology, the attempt to argue from the phenomena or good qualities of nature to God.
Barth rejected this entirely, because it sinks inevitably into idolatry, describing God in our own
terms. Liberalism had done something similar, keeping only the biblical teaching that suited its
refined idea of God, and the outcome was the same as in Nazi theology: God not only likes us, he is
like us. Shelley sees the Nazi genocide as a demonic form of idolatry, which its believers took for
a divinely approved cleansing. Aberrant theology has consequences.

The Confessing Church planned no campaign of resistance. It was aimed at the heresies of the German
Christians, and its leaders repeatedly affirmed loyalty to the state and congratulated Hitler on his
political moves. Following the Lutheran habit of supporting the ruling power, it did not set itself
up as a rival church but as a body defending orthodox faith against innovation. Harassed by the
Gestapo and disowned by most Protestant leaders, it led a perilous life; in 1935 seven hundred of its
pastors were arrested. Yet its existence embarrassed the Nazis, and its witness to Christ's lordship
over the world was an implicit challenge to totalitarianism. When M\"uller failed to unite the
Protestants, Hitler turned increasingly to the anti-Christian Nazis who called Nazism the true
fulfillment of Christianity. In 1935 a Ministry of Church Affairs was created under the lawyer Hanns
Kerrl, who declared, ``National Socialism is the doing of God's will. God's will reveals itself in
German blood. True Christianity is represented by the party.''

Catholic leaders, alarmed by the ``new heathenism'' and by growing restrictions, turned to the
Vatican. In 1937 Pius XI, with the help of his secretary of state Cardinal Eugenio Pacelli, soon to
be Pius XII, drafted the encyclical \emph{Mit brennender Sorge} (``With deep anxiety''), the first
major church document to criticize Nazism. Smuggled into Germany, it was read from every Catholic
pulpit on Palm Sunday before a single copy had fallen into Nazi hands. As Richard Pierard explains,
it protested the oppression of the church and called on Catholics to resist the idolatrous cult of
race and state, to stand against the perversion of Christian doctrine and morality, and to remain
loyal to Christ, his church and Rome. Hitler was furious at first, then chose to avoid a break with
Rome by meeting the encyclical with silence. Sure of the support of the Catholic laity, he simply
increased the pressure on the churches so that organized resistance became impossible.

\ssection{The Russian Revolution}
In the same decades the Bolsheviks built a totalitarian system of the left. It shared many traits
with Hitler's Germany: dictatorial leadership, a single centralized party, ruthless terror,
propaganda, censorship, a controlled economy and hostility to all organized religion. Its ideology
differed in stressing the working class, revolution as the means of social change and the ideal of a
classless society. From 1917 to 1924 its mastermind was Vladimir Ilich Ulyanov (1870--1924), known as
Lenin, an exiled socialist who returned from Switzerland in 1917. He believed in violence as the
chief weapon of Marxism, and to destroy capitalism he created a highly disciplined Communist party,
first to overthrow the old institutions and then to run the new government. Party and state became
one. After his death Leon Trotsky and Joseph Stalin fought for the succession, and by 1927 Stalin was
unrivaled dictator. Determined to make the Soviet Union an industrial nation able to compete with the
capitalist countries, he proved even more ruthless than Hitler, using secret police and labor camps
to crush the slightest dissent and remove every possible rival.

Communism confronted Christianity with a worldview founded on atheism. In Marxist-Leninist theory
religion is false consciousness, an illusory reflection of a world divided into classes, which will
die a natural death once society reaches its ``normal'' state. Yet Communists fight religion
actively, because the party regards itself as the embodiment of Marxism and can allow no part of
society outside its control. Institutional religion is a reactionary force and must be smashed. The
Russian Orthodox Church had for centuries been the state church, with the tsar in theory its head.
The Bolsheviks confiscated church lands, ended state subsidies, decreed civil marriage and forbade
religious education of the young. How could a church survive that could not teach its own children
amid atheist indoctrination? The newly elected Patriarch Tikhon, in his first pastoral letter,
declared war on the state, and riots followed in almost every city. In the unorganized violence of
the first six years twenty-eight bishops and more than a thousand priests were killed.

The Law on Religious Associations of 1929 set strict limits on church activity and in effect denied
the churches any influence on society. Years of intense persecution followed, and thousands of
clergy were imprisoned or killed during the collectivization of agriculture and Stalin's purges. The
constitution of 1936 restored the clergy's right to vote, but the ``servants of religion'' remained
second-class citizens, accused of exploiting the ignorance of the toilers and hounded by the secret
police as ``clerico-fascists.'' By 1939 propaganda, anti-religious laws and terror had brought the
Orthodox Church to the brink of disintegration; the Lutherans were almost wiped out, and the Baptists
and Evangelical Christians were ravaged. Revulsion grew in the West. In March 1937 Pius XI issued
the encyclical \emph{Divini Redemptoris} against the ``errors of communism.'' He expressed sympathy
for the Russian people, offered Catholic teaching as the alternative, and, against Stalin's
``popular front'' policy, declared that communism is ``intrinsically wrong'' and that no one who
would save Christian civilization may collaborate with it in anything. Coming within days of the
encyclical on Nazi Germany, this set the Vatican firmly on the side of persecuted believers in the
totalitarian countries.

\ssection{The Impact of World War II}
The war began in 1939 with the German invasion of Poland. Hitler had made common cause with
Mussolini's fascist Italy and the militarist clique ruling Japan, all three bent on expansion. He
signed a nonaggression pact with Stalin just before the invasion and within two years turned on him,
which forced Communist Russia into alliance with France, Britain and the United States. As in the
first war, Christians on both sides took part, though with less crusading zeal. German church
officials stayed conciliatory and did nothing to relieve the suffering of Christians. Hitler's
closest advisers, Bormann, Himmler and Heydrich, worked toward a ``final settlement'' of the church
question: the churches subordinated to the new order, the clergy stripped of privileges, and
Christianity left to die what Hitler called ``a natural death.'' In occupied eastern Europe priests,
pastors and devout laymen were treated as common criminals, and thousands were executed or sent to
camps. Only the needs of war and of popular support kept the Nazis from rooting out religion in
Germany itself. The German churches' resistance was amazingly meager. Their concern with individual
personal faith, their traditional submission to the state, and a conservatism that rejected every
left-wing proposal for reform let them accept the Nazi claim to be the only alternative to
communism.

Wartime Russia was a striking contrast. Stalin saw what the churches could do for morale, for
absorbing the territories gained in the war and for later foreign policy. He allowed them to rebuild
their organizations, collect funds and give some private religious instruction to children. In 1943
the government created a Council for Affairs of the Russian Orthodox Church, and a year later a
similar council for other groups, to maintain relations, draft regulations and see the laws
enforced. In 1945 the Orthodox Church and other bodies regained legal status, with the right to hold
property and make liturgical objects. Still closely supervised, the Orthodox Church was better off
than at any time since the Revolution.

On the whole the war was devastating for Christianity, physically and morally. Thousands of churches
were destroyed, clergy were killed, and believers were persecuted or uprooted. Violence rose with
armored vehicles, fire-bombing, guided missiles and the atomic bomb, and millions died. Shelley
rejects the view, popular in American universities, that the Soviet regime was a corrupt but morally
equivalent counterpart of the United States, a flawed socialism beside an imperfect capitalism.
Communist killings, still being counted, exceed the Nazi toll of six million, he says, by a factor
of ten or more. The carnage led many to ask whether a just war is still possible and to suggest that
Christian endorsement of war only intensifies it. Some Christians joined in efforts to bring nations
and churches together, such as the United Nations (1945) and the World Council of Churches (1948),
but national passions frustrated them.

After the atomic bomb brought Japan to terms and the Allied and Communist armies toppled the Third
Reich, the victors entered the ``Cold War.'' The United States, leading the Western democracies,
adopted the policy of containment and was ready to resist communism almost anywhere, by force if
necessary, until the hydrogen bomb and long-range missiles made war between the superpowers costly
beyond reckoning. In Richard Pierard's summary, a rivalry of great powers quickly became a contest of
ideologies, and each side had the support of its Christians. The Soviets preached communism and the
liberation of peoples oppressed by ``imperialists'' with messianic fervor. Russian Orthodox leaders
sought to bring churches in other countries under their jurisdiction, repeated the Moscow line
abroad, praised conditions at home and were active in Soviet-sponsored peace campaigns; with the
stress on ``peaceful coexistence'' after Stalin's death in 1953 they took part in ecumenical affairs
and in 1961 joined the World Council of Churches. The West answered with the counter-ideology of
anticommunism. It assumed a universal conspiracy centered in Moscow behind all revolutionary unrest,
and it was above all an American response to the stalemate, born of anxious frustration at America's
inability to spread liberal democracy to every nation. So the wall remained.

\begin{keyfacts}
\keypeople{Arnold Toynbee; Woodrow Wilson; Adolf Hitler; Ludwig M\"uller; Martin Niem\"oller; Karl
Barth; Hanns Kerrl; Pius XI; Lenin; Joseph Stalin; Patriarch Tikhon; Nikita Khrushchev.}
\keydates{%
  \kd{1914}{Outbreak of the First World War}
  \kd{1917}{Bolshevik Revolution in Russia; the United States enters the war}
  \kd{1929}{Soviet Law on Religious Associations}
  \kd{1933}{Hitler becomes chancellor; concordat with the Vatican}
  \kd{1934}{Barmen Declaration of the Confessing Church}
  \kd{1937}{Pius XI's encyclicals against Nazism and communism}
  \kd{1939--1945}{Second World War; the Holocaust}
  \kd{1953}{Death of Stalin}
  \kd{1961}{Berlin Wall built; Russian Orthodox Church joins the World Council of Churches}}
\keyterms{Ideology; secularization; totalitarianism; nationalism; fascism; Nazism; leadership
principle; anti-Semitism; Holocaust; German Christians; Confessing Church; Barmen Declaration;
natural theology; Bolsheviks; false consciousness; Cold War; containment; anticommunism.}
\end{keyfacts}
